% !TEX program = xelatex
% !TEX root = mainsempro.tex

% --- Versi Seminar Proposal (Bab 1–3 saja) ---

% !TEX root = main.tex

% =============================
\documentclass[12pt,a4paper]{report}

% --- Font (Times New Roman asli) ---
\usepackage{fontspec}
\setmainfont{Times New Roman}[Ligatures=TeX,Renderer=Basic]

\usepackage[american,indonesian]{babel}

\usepackage{csquotes}
\usepackage[style=apa,backend=biber]{biblatex}
\addbibresource{references.bib}

% --- MAP bahasa indonesian ke mapping 'american-apa' yang tersedia ---
\DeclareLanguageMapping{indonesian}{american-apa}

% --- Override label bibliografi ke Bahasa Indonesia ---
\DefineBibliographyStrings{indonesian}{%
  bibliography = {DAFTAR PUSTAKA},
  references = {DAFTAR PUSTAKA},
  and          = {\&},
}

% Pilih bahasa tampilan dokumen jadi indonesian
\AtBeginDocument{\selectlanguage{indonesian}}

% --- Margin ---
\usepackage{geometry}
\geometry{
  a4paper,
  left=4cm,
  right=3cm,
  top=4cm,
  bottom=3cm
}

% --- Spasi 1.5 ---
\usepackage{setspace}
\onehalfspacing

% --- Paragraf (indentasi) ---
\usepackage{indentfirst}
\setlength{\parindent}{1cm}
\setlength{\parskip}{0pt}

% --- Matematika & fisika ---
\usepackage{amsmath,amssymb,amsfonts,mathtools}
\usepackage{physics}
\usepackage{bm}
\usepackage{tensor}
\usepackage{siunitx}

% --- Gambar & grafik ---
\usepackage{graphicx}
\usepackage{caption}
\usepackage{subcaption}
\usepackage{float}

% --- Tabel ---
\usepackage{booktabs}
\usepackage{multirow}
\usepackage{array}

% --- Listing kode Python ---
\usepackage{listings}
\usepackage{xcolor}

\definecolor{codebg}{RGB}{245,245,245}
\definecolor{codegreen}{RGB}{0,128,0}
\definecolor{codegray}{RGB}{128,128,128}
\definecolor{codepurple}{RGB}{128,0,128}

\lstdefinestyle{pythonstyle}{
  backgroundcolor=\color{codebg},
  basicstyle=\ttfamily\small,
  breaklines=true,
  captionpos=b,
  commentstyle=\color{codegreen},
  keywordstyle=\color{codepurple}\bfseries,
  stringstyle=\color{red!70!black},
  numberstyle=\tiny\color{codegray},
  numbers=left,
  numbersep=5pt,
  frame=single,
  rulecolor=\color{black!30},
  tabsize=4,
  language=Python,
  showstringspaces=false,
}
\lstset{style=pythonstyle}

% --- Hyperlink & referensi ---
\usepackage[hidelinks]{hyperref}
\usepackage{cleveref}

% --- Penomoran persamaan per bab ---
\numberwithin{equation}{chapter}

% --- Enumerate custom ---
\usepackage{enumitem}

% --- Format judul bab ---
\usepackage{titlesec}

\renewcommand{\thechapter}{\arabic{chapter}}
\renewcommand{\thesection}{\thechapter.\arabic{section}}
\renewcommand{\thesubsection}{\thesection.\arabic{subsection}}

\titleformat{\chapter}[display]
  {\normalfont\centering\bfseries\normalsize}
  {\MakeUppercase{BAB \Roman{chapter}}}
  {0pt}
  {\vspace{1ex}\MakeUppercase}
\titlespacing*{\chapter}{0pt}{0pt}{20pt}

\titleformat{\section}
  {\normalfont\bfseries\normalsize}
  {\thesection}{1em}{}
\titlespacing*{\section}{0pt}{12pt}{6pt}

\titleformat{\subsection}
  {\normalfont\bfseries\normalsize}
  {\thesubsection}{1em}{}
\titlespacing*{\subsection}{0pt}{10pt}{4pt}

\usepackage[none]{hyphenat}
\sloppy


% --- Informasi skripsi ---
\title{\textbf{Kajian Analitik dan Penyelesaian Numerik Efek Lense-Thirring pada Metrik Kerr: Studi Kasus Gravity Probe B}}
\author{Pramudya Jesril Pratama}
\date{\today}

\begin{document}

% ---------- Judul (titlepage) ----------
\maketitle

% ---------- Penomoran HALAMAN MUKA: Roman kecil ----------
\pagenumbering{roman}

% !TEX root = ../main.tex

\chapter*{Abstrak}
\addcontentsline{toc}{chapter}{Abstrak}

Relativitas Umum memprediksi bahwa massa berotasi menginduksi penyeretan kerangka acuan inersia lokal, yang dikenal sebagai efek Lense-Thirring atau \textit{frame-dragging}.
Efek ini telah dikonfirmasi secara eksperimental oleh misi \textit{Gravity Probe B} (GP-B) dengan hasil pengukuran sebesar $37{,}2 \pm 7{,}2$ milliarcsecond per tahun (mas/yr) untuk presesi \textit{frame-dragging} dan $6602 \pm 18$ mas/yr untuk presesi geodesik.
Penelitian ini menyajikan kajian analitik terhadap persamaan geodesik dan transportasi paralel tensor spin pada metrik Kerr, serta penyelesaian numeriknya untuk konfigurasi orbit Bumi.
Metrik Kerr dalam koordinat Boyer-Lindquist diturunkan secara eksplisit dalam bentuk matriks $4 \times 4$, dan simbol-simbol Christoffel yang relevan dihitung untuk memperoleh persamaan gerak.
Lagrangian metrik Kerr digunakan untuk menurunkan persamaan geodesik melalui metode Euler-Lagrange, yang kemudian direduksi menjadi sistem persamaan diferensial biasa orde satu.
Aproksimasi medan lemah diterapkan untuk memperoleh ekspresi analitik laju presesi Lense-Thirring $\Omega_{\text{LT}} \approx 2GJ/(c^2 r^3)$.
Sistem persamaan diferensial yang dihasilkan diselesaikan secara numerik menggunakan metode Runge-Kutta orde empat yang diimplementasikan dalam Python.
Hasil simulasi menunjukkan laju presesi geodesik sebesar $\sim\!6600$ mas/yr dan \textit{frame-dragging} sebesar $\sim\!37$ mas/yr, konsisten dengan data eksperimen GP-B dalam batas galat pengukuran.
Analisis kestabilan numerik dan sensitivitas terhadap parameter rotasi Bumi turut disajikan untuk memvalidasi keandalan metode komputasi yang digunakan.

\vspace{1em}
\noindent\textbf{Kata Kunci:} Relativitas Umum, Metrik Kerr, Efek Lense-Thirring, \textit{Frame-Dragging}, Presesi Geodesik, Gravity Probe B, Penyelesaian Numerik.


\tableofcontents
\listoffigures
\listoftables

% ---------- Isi utama: penomoran Arab ----------
\cleardoublepage
\pagenumbering{arabic}
\setcounter{page}{1}

% === BAB I: PENDAHULUAN ===
% !TEX root = ../main.tex

\chapter{Pendahuluan}
\label{chap:pendahuluan}

% ============================================================
\section{Latar Belakang}
\label{sec:latar-belakang}

Teori Relativitas Umum (\textit{General Relativity}, GR) yang dirumuskan oleh Einstein pada tahun 1915 merupakan teori gravitasi yang mendeskripsikan interaksi gravitasi sebagai manifestasi kelengkungan ruang-waktu empat dimensi. Dalam kerangka ini, distribusi massa dan energi menentukan geometri ruang-waktu melalui persamaan medan Einstein $G_{\mu\nu} = \frac{8\pi G}{c^4}\, T_{\mu\nu}$, di mana $G_{\mu\nu}$ adalah tensor Einstein dan $T_{\mu\nu}$ adalah tensor energi-momentum \parencite{Einstein1916GR, Misner1973Gravitation}. Gerak benda uji dalam ruang-waktu melengkung ditentukan oleh persamaan geodesik yang bergantung pada simbol Christoffel $\Gamma^\mu_{\nu\rho}$, yang diturunkan dari tensor metrik $g_{\mu\nu}$ \parencite{Carroll2004SpacetimeGeometry}.

Salah satu konsekuensi penting dari GR adalah prediksi bahwa massa berotasi tidak hanya melengkungkan ruang-waktu, tetapi juga menyeret kerangka acuan inersia lokal di sekitarnya. Efek ini dikenal sebagai \textit{frame-dragging} atau efek Lense-Thirring, yang pertama kali diprediksi secara teoretis oleh \textcite{Lense1918LenseThirring}. Secara matematis, efek ini muncul dari suku campuran $g_{t\phi}$ pada metrik Kerr yang mendeskripsikan ruang-waktu di sekitar massa berotasi \parencite{Kerr1963KerrMetric}. Selain efek Lense-Thirring, GR juga memprediksi adanya presesi geodesik (\textit{geodetic precession} atau efek de Sitter), yaitu perubahan orientasi giroskop akibat geraknya melalui ruang-waktu yang melengkung oleh massa statis \parencite{Carroll2004SpacetimeGeometry, Will2014Confrontation}.

Kedua efek ini telah diverifikasi secara eksperimental melalui misi \textit{Gravity Probe B} (GP-B) yang diluncurkan oleh NASA pada tahun 2004. Misi ini menggunakan empat giroskop kuarsa berpresisi tinggi dalam orbit polar di sekitar Bumi pada ketinggian $\sim$642~km. Hasil pengukuran GP-B menunjukkan laju presesi geodesik sebesar $6602 \pm 18$ milliarcsecond per tahun (mas/yr) dan laju presesi \textit{frame-dragging} sebesar $37{,}2 \pm 7{,}2$~mas/yr, keduanya konsisten dengan prediksi GR \parencite{Everitt2011GravityProbeB}.

Meskipun hasil eksperimen GP-B telah dikonfirmasi, kajian analitik terhadap proses derivasi matematis yang mendasari prediksi tersebut tetap memiliki nilai penting. Prediksi laju presesi Lense-Thirring memerlukan serangkaian langkah derivasi yang meliputi: (i)~penulisan eksplisit tensor metrik Kerr dalam koordinat Boyer-Lindquist, (ii)~perhitungan simbol Christoffel dari komponen-komponen metrik, (iii)~penurunan persamaan geodesik melalui metode Lagrangian, (iv)~formulasi persamaan transportasi paralel tensor spin, dan (v)~reduksi ke aproksimasi medan lemah untuk memperoleh ekspresi analitik laju presesi. Setiap langkah tersebut melibatkan operasi kalkulus tensor yang tidak trivial dan jarang disajikan secara lengkap dalam satu kajian terpadu \parencite{Carroll2004SpacetimeGeometry, Schutz2009FirstCourse}.

Di sisi lain, persamaan diferensial yang dihasilkan dari derivasi tersebut bersifat sangat nonlinier, sehingga penyelesaian analitik eksak hanya dimungkinkan pada kasus-kasus tertentu dengan simetri tinggi. Untuk konfigurasi orbit yang realistis, diperlukan penyelesaian numerik menggunakan metode seperti Runge-Kutta orde empat atau integrator simplektik \parencite{Butcher2016NumericalMethods, Hairer2006GeometricIntegration}. Implementasi numerik ini dapat dilakukan dengan bantuan pustaka komputasi Python seperti \texttt{EinsteinPy} dan \texttt{SciPy} \parencite{Bapat2020EinsteinPy, SciPy2020}.

Berdasarkan latar belakang tersebut, penelitian ini menyajikan kajian analitik terhadap persamaan geodesik dan transportasi paralel tensor spin pada metrik Kerr, serta penyelesaian numerik sistem persamaan diferensial yang dihasilkan untuk konfigurasi orbit Bumi. Hasil komputasi numerik kemudian dibandingkan dengan data eksperimen GP-B sebagai validasi konsistensi antara derivasi analitik dan pengamatan eksperimental.

% ============================================================
\section{Batasan Masalah}
\label{sec:batasan-masalah}

Penelitian ini dibatasi pada hal-hal berikut:
\begin{enumerate}[label=\arabic*)]
    \item Kajian dilakukan dalam kerangka Relativitas Umum dengan menggunakan metrik Kerr sebagai representasi ruang-waktu di sekitar Bumi yang berotasi. Metrik Schwarzschild digunakan sebagai kasus batas non-rotasi ($a = 0$) untuk keperluan analisis sensitivitas.
    \item Model ruang-waktu dibatasi pada orbit di sekitar Bumi dengan mengabaikan pengaruh benda langit lain.
    \item Analisis difokuskan pada dua efek relativistik utama: presesi geodesik dan efek Lense-Thirring (\textit{frame-dragging}).
    \item Derivasi analitik dilakukan dalam kerangka kalkulus tensor. Penyelesaian numerik menggunakan metode Runge-Kutta yang diimplementasikan dalam Python.
    \item Penelitian bersifat kajian analitik-numerik, tanpa melibatkan pengamatan eksperimental langsung.
\end{enumerate}

% ============================================================
\section{Rumusan Masalah}
\label{sec:rumusan-masalah}

Berdasarkan latar belakang dan batasan masalah yang telah diuraikan, rumusan masalah dalam penelitian ini adalah sebagai berikut:
\begin{enumerate}[label=\arabic*)]
    \item Bagaimana menurunkan persamaan geodesik dan persamaan transportasi paralel tensor spin dari metrik Kerr melalui metode Lagrangian dan kalkulus tensor?
    \item Bagaimana mereduksi metrik Kerr ke dalam aproksimasi medan lemah untuk memperoleh ekspresi analitik laju presesi Lense-Thirring pada kasus Bumi?
    \item Bagaimana menyelesaikan sistem persamaan diferensial geodesik dan transportasi paralel spin secara numerik, serta memvalidasi hasilnya terhadap data eksperimen \textit{Gravity Probe B}?
\end{enumerate}

% ============================================================
\section{Tujuan Penelitian}
\label{sec:tujuan-penelitian}

Tujuan dari penelitian ini adalah sebagai berikut:
\begin{enumerate}[label=\arabic*)]
    \item Menurunkan persamaan geodesik dari Lagrangian metrik Kerr melalui persamaan Euler-Lagrange, serta merumuskan persamaan transportasi paralel tensor spin dalam bentuk matriks.
    \item Melakukan reduksi metrik Kerr ke aproksimasi medan lemah dan menurunkan ekspresi analitik laju presesi Lense-Thirring $\Omega_{\text{LT}} \approx 2GJ/(c^2 r^3)$ untuk konfigurasi Bumi.
    \item Menyelesaikan sistem persamaan diferensial yang diperoleh secara numerik menggunakan metode Runge-Kutta, dan membandingkan hasil komputasi dengan data pengukuran \textit{Gravity Probe B}.
\end{enumerate}

% ============================================================
\section{Manfaat Penelitian}
\label{sec:manfaat-penelitian}

Manfaat dari penelitian ini adalah sebagai berikut:
\begin{enumerate}[label=\arabic*)]
    \item Bagi ilmu pengetahuan: menyediakan kajian terpadu yang menjabarkan derivasi analitik efek Lense-Thirring dari metrik Kerr secara langkah demi langkah, mulai dari tensor metrik hingga ekspresi laju presesi, disertai verifikasi numerik.
    \item Bagi penulis: mengembangkan pemahaman mendalam di bidang kalkulus tensor, Relativitas Umum, dan metode penyelesaian numerik persamaan diferensial.
    \item Bagi akademis: menghasilkan dokumen referensi yang dapat digunakan sebagai bahan pembelajaran mengenai derivasi efek relativistik pada metrik Kerr dan implementasi komputasinya.
\end{enumerate}


% === BAB II: TINJAUAN PUSTAKA ===
% !TEX root = ../main.tex

\chapter{Tinjauan Pustaka}
\label{chap:tinjauan}

% ============================================================
\section{Relativitas Umum Einstein}
\label{sec:relativitas-umum}

Relativitas Umum (GR) mendeskripsikan gravitasi sebagai kelengkungan ruang-waktu empat dimensi yang ditentukan oleh distribusi massa dan energi. Hubungan antara geometri ruang-waktu dan sumber gravitasi dinyatakan oleh persamaan medan Einstein \parencite{Einstein1916GR}:
\begin{equation}
    R_{\mu\nu} - \frac{1}{2}\, R\, g_{\mu\nu} = \frac{8\pi G}{c^4}\, T_{\mu\nu},
    \label{eq:efe}
\end{equation}
di mana $R_{\mu\nu}$ adalah tensor Ricci, $R = g^{\mu\nu}R_{\mu\nu}$ adalah skalar kelengkungan, $g_{\mu\nu}$ adalah tensor metrik, dan $T_{\mu\nu}$ adalah tensor energi-momentum. Tensor Ricci diperoleh dari kontraksi tensor Riemann:
\begin{equation}
    R^\rho_{\ \sigma\mu\nu} = \partial_\mu \Gamma^\rho_{\nu\sigma}
    - \partial_\nu \Gamma^\rho_{\mu\sigma}
    + \Gamma^\rho_{\mu\lambda}\Gamma^\lambda_{\nu\sigma}
    - \Gamma^\rho_{\nu\lambda}\Gamma^\lambda_{\mu\sigma},
    \label{eq:riemann}
\end{equation}
dengan $R_{\mu\nu} = R^\rho_{\ \mu\rho\nu}$ \parencite{Carroll2004SpacetimeGeometry, Misner1973Gravitation}.

Gerak partikel bebas dalam ruang-waktu melengkung mengikuti lintasan geodesik, yaitu kurva yang memenuhi:
\begin{equation}
    \frac{d^2 x^\mu}{d\tau^2}
    + \Gamma^\mu_{\nu\rho}\,\frac{dx^\nu}{d\tau}\,\frac{dx^\rho}{d\tau} = 0,
    \label{eq:geodesic}
\end{equation}
di mana $\tau$ adalah waktu proper dan $\Gamma^\mu_{\nu\rho}$ adalah simbol Christoffel yang didefinisikan pada Persamaan~\eqref{eq:christoffel-def} \parencite{Schutz2009FirstCourse}.


% ============================================================
\section{Tensor Metrik dan Simbol Christoffel}
\label{sec:metrik-christoffel}

\subsection{Simbol Christoffel}
\label{subsec:christoffel-umum}

Simbol Christoffel jenis kedua $\Gamma^\mu_{\nu\rho}$ menggambarkan koneksi Levi-Civita pada manifold pseudo-Riemannian dan didefinisikan sebagai:
\begin{equation}
    \Gamma^\mu_{\nu\rho} = \frac{1}{2}\, g^{\mu\sigma}
    \left(
        \frac{\partial g_{\sigma\nu}}{\partial x^\rho}
      + \frac{\partial g_{\sigma\rho}}{\partial x^\nu}
      - \frac{\partial g_{\nu\rho}}{\partial x^\sigma}
    \right).
    \label{eq:christoffel-def}
\end{equation}
Simbol ini bersifat simetris terhadap dua indeks bawah, $\Gamma^\mu_{\nu\rho} = \Gamma^\mu_{\rho\nu}$, dan berjumlah paling banyak 40 komponen independen dalam ruang-waktu empat dimensi. Untuk metrik diagonal, banyak komponen bernilai nol, namun metrik Kerr memiliki suku campuran $g_{t\phi} \neq 0$ yang menghasilkan komponen Christoffel tambahan yang tidak trivial \parencite{Carroll2004SpacetimeGeometry}.


\subsection{Metrik Schwarzschild}
\label{subsec:metrik-schwarzschild}

Solusi Schwarzschild mendeskripsikan ruang-waktu di sekitar massa statis tak berotasi. Dalam koordinat bola $(t, r, \theta, \phi)$, elemen garisnya diberikan oleh \parencite{Schwarzschild1916Field}:
\begin{equation}
    ds^2 = \left(1 - \frac{r_s}{r}\right) c^2\,dt^2
         - \left(1 - \frac{r_s}{r}\right)^{-1} dr^2
         - r^2\, d\theta^2
         - r^2 \sin^2\!\theta\, d\phi^2,
    \label{eq:schwarzschild}
\end{equation}
dengan $r_s = 2GM/c^2$ adalah radius Schwarzschild. Metrik ini bersifat diagonal dan menghasilkan efek presesi geodesik, namun tidak mengandung efek \textit{frame-dragging} karena tidak adanya suku campuran $dt\,d\phi$.


\subsection{Metrik Kerr}
\label{subsec:metrik-kerr}

Metrik Kerr merupakan solusi eksak persamaan medan Einstein untuk massa berotasi tanpa muatan listrik, ditemukan oleh \textcite{Kerr1963KerrMetric}. Dalam koordinat Boyer-Lindquist $(t, r, \theta, \phi)$, elemen garisnya dapat dituliskan sebagai:
\begin{equation}
\begin{split}
    ds^2 &= \left(1 - \frac{r_s r}{\Sigma}\right) c^2\,dt^2
          + \frac{2\,r_s\, r\, a \sin^2\!\theta}{\Sigma}\; c\,dt\,d\phi
          - \frac{\Sigma}{\Delta}\, dr^2
          - \Sigma\, d\theta^2 \\
         &\quad - \left(r^2 + a^2 + \frac{r_s\, r\, a^2 \sin^2\!\theta}{\Sigma}\right)
           \sin^2\!\theta\, d\phi^2,
\end{split}
    \label{eq:kerr-line}
\end{equation}
dengan fungsi-fungsi bantu:
\begin{equation}
    \Sigma \equiv r^2 + a^2 \cos^2\!\theta, \qquad
    \Delta \equiv r^2 - r_s\, r + a^2,
    \label{eq:sigma-delta}
\end{equation}
di mana $r_s = 2GM/c^2$ dan $a = J/(Mc)$ adalah parameter rotasi spesifik, dengan $J$ menyatakan momentum sudut total massa pusat \parencite{Kerr1963KerrMetric, Carroll2004SpacetimeGeometry}.

Dengan menggunakan konvensi koordinat $x^\mu = (x^0, x^1, x^2, x^3) = (ct, r, \theta, \phi)$, tensor metrik Kerr $g_{\mu\nu}$ dapat dituliskan dalam bentuk matriks $4 \times 4$ sebagai:
\begin{equation}
    g_{\mu\nu} =
    \begin{pmatrix}
        1 - \dfrac{r_s r}{\Sigma} & 0 & 0 & \dfrac{r_s\, r\, a \sin^2\!\theta}{\Sigma} \\[12pt]
        0 & -\dfrac{\Sigma}{\Delta} & 0 & 0 \\[12pt]
        0 & 0 & -\Sigma & 0 \\[12pt]
        \dfrac{r_s\, r\, a \sin^2\!\theta}{\Sigma} & 0 & 0
        & -\!\left(r^2 + a^2 + \dfrac{r_s\, r\, a^2 \sin^2\!\theta}{\Sigma}\right)\sin^2\!\theta
    \end{pmatrix}.
    \label{eq:kerr-matrix}
\end{equation}

Beberapa sifat penting dari matriks metrik ini:
\begin{enumerate}[label=(\roman*)]
    \item Komponen $g_{t\phi} = g_{\phi t} = r_s\, r\, a \sin^2\!\theta\,/\,\Sigma$ tidak bernilai nol jika $a \neq 0$. Suku campuran inilah yang bertanggung jawab atas efek \textit{frame-dragging}.
    \item Pada limit $a \to 0$, matriks di atas tereduksi menjadi metrik Schwarzschild diagonal (Persamaan~\ref{eq:schwarzschild}).
    \item Komponen $g_{rr}$ memiliki singularitas koordinat pada $\Delta = 0$, yang terjadi pada horizon peristiwa.
\end{enumerate}

Metrik invers $g^{\mu\nu}$ diperlukan untuk menghitung simbol Christoffel. Dengan memanfaatkan struktur blok-diagonal ($2 \times 2$ pada sub-ruang $t$-$\phi$ dan diagonal pada $r$-$\theta$), komponen-komponen metrik invers diperoleh sebagai:
\begin{equation}
\begin{aligned}
    g^{tt} &= \frac{1}{\Delta}\left(r^2 + a^2 + \frac{r_s\, r\, a^2 \sin^2\!\theta}{\Sigma}\right), &\qquad
    g^{t\phi} &= \frac{r_s\, r\, a}{\Sigma\,\Delta}, \\[6pt]
    g^{rr} &= -\frac{\Delta}{\Sigma}, &\qquad
    g^{\theta\theta} &= -\frac{1}{\Sigma}, \\[6pt]
    g^{\phi\phi} &= -\frac{1}{\Delta\sin^2\!\theta}\left(1 - \frac{r_s\, r}{\Sigma}\right), &\qquad
    & \text{komponen lain} = 0.
\end{aligned}
    \label{eq:kerr-inverse}
\end{equation}


\subsection{Contoh Perhitungan Simbol Christoffel Metrik Kerr}
\label{subsec:christoffel-kerr-contoh}

Perhitungan simbol Christoffel untuk metrik Kerr menghasilkan sejumlah besar komponen yang tidak bernilai nol. Pada bagian ini disajikan derivasi eksplisit untuk tiga komponen representatif yang relevan dengan efek \textit{frame-dragging}.

\paragraph{Contoh 1: $\Gamma^t_{tr}$.}
Menggunakan Persamaan~\eqref{eq:christoffel-def} dengan $\mu = t$, $\nu = t$, $\rho = r$:
\begin{equation}
    \Gamma^t_{tr} = \frac{1}{2}\,g^{t\sigma}
    \left(\partial_r g_{\sigma t} + \partial_t g_{\sigma r} - \partial_\sigma g_{tr}\right).
\end{equation}
Karena metrik tidak bergantung pada $t$ (stasioner), suku $\partial_t g_{\sigma r} = 0$. Kontribusi non-nol berasal dari $\sigma = t$ dan $\sigma = \phi$:
\begin{equation}
    \Gamma^t_{tr} = \frac{1}{2}\,g^{tt}\,\partial_r g_{tt}
                  + \frac{1}{2}\,g^{t\phi}\,\partial_r g_{\phi t}.
    \label{eq:christoffel-ttr}
\end{equation}
Turunan parsial $\partial_r g_{tt}$ dan $\partial_r g_{t\phi}$ dapat dihitung dari Persamaan~\eqref{eq:kerr-matrix}. Sebagai contoh:
\begin{equation}
    \frac{\partial g_{tt}}{\partial r}
    = \frac{\partial}{\partial r}\left(1 - \frac{r_s\, r}{\Sigma}\right)
    = -r_s\,\frac{\Sigma - r(2r)}{\Sigma^2}
    = \frac{r_s(r^2 - a^2\cos^2\!\theta)}{\Sigma^2},
    \label{eq:dgtt-dr}
\end{equation}
mengingat $\partial_r \Sigma = 2r$. Substitusi ke Persamaan~\eqref{eq:christoffel-ttr} bersama dengan ekspresi $g^{tt}$ dan $g^{t\phi}$ dari Persamaan~\eqref{eq:kerr-inverse} menghasilkan:
\begin{equation}
    \Gamma^t_{tr} = \frac{r_s(r^2 - a^2\cos^2\!\theta)(r^2 + a^2)}{2\,\Sigma^2\,\Delta}
    + \frac{r_s^2\, r\, a^2\sin^2\!\theta\,(r^2 - a^2\cos^2\!\theta)}{2\,\Sigma^3\,\Delta}.
    \label{eq:gamma-ttr-result}
\end{equation}

\paragraph{Contoh 2: $\Gamma^\phi_{tr}$.}
Komponen ini menggambarkan kopling langsung antara gerak radial dan gerak azimut akibat rotasi sumber:
\begin{equation}
    \Gamma^\phi_{tr} = \frac{1}{2}\,g^{\phi\sigma}
    \left(\partial_r g_{\sigma t} + \partial_t g_{\sigma r} - \partial_\sigma g_{tr}\right).
\end{equation}
Dengan $\partial_t = 0$ dan kontribusi dari $\sigma = t$ serta $\sigma = \phi$:
\begin{equation}
    \Gamma^\phi_{tr} = \frac{1}{2}\,g^{\phi t}\,\partial_r g_{tt}
                     + \frac{1}{2}\,g^{\phi\phi}\,\partial_r g_{\phi t}.
    \label{eq:christoffel-phitr}
\end{equation}
Perhatikan bahwa $g^{\phi t} = g^{t\phi}$. Substitusi komponen metrik invers (Persamaan~\ref{eq:kerr-inverse}) dan turunan yang sesuai memberikan ekspresi dalam $r$, $\theta$, $\Sigma$, dan $\Delta$. Komponen ini bernilai nol pada limit $a \to 0$ (metrik Schwarzschild), yang menegaskan bahwa kopling $r$-$\phi$ murni berasal dari rotasi massa.

\paragraph{Contoh 3: $\Gamma^t_{\theta\phi}$.}
Komponen ini menunjukkan kopling antara gerak polar dan azimut melalui koordinat waktu:
\begin{equation}
    \Gamma^t_{\theta\phi} = \frac{1}{2}\,g^{t\sigma}
    \left(\partial_\phi g_{\sigma\theta} + \partial_\theta g_{\sigma\phi} - \partial_\sigma g_{\theta\phi}\right).
\end{equation}
Karena metrik tidak bergantung pada $\phi$ (simetri aksial) dan $g_{r\theta} = g_{r\phi} = 0$, maka:
\begin{equation}
    \Gamma^t_{\theta\phi} = \frac{1}{2}\,g^{tt}\,\partial_\theta g_{t\phi}
                          + \frac{1}{2}\,g^{t\phi}\,\partial_\theta g_{\phi\phi}.
    \label{eq:christoffel-tthetaphi}
\end{equation}
Turunan $\partial_\theta g_{t\phi}$ memerlukan:
\begin{equation}
    \frac{\partial g_{t\phi}}{\partial\theta}
    = \frac{\partial}{\partial\theta}
    \!\left(\frac{r_s\, r\, a\sin^2\!\theta}{\Sigma}\right)
    = \frac{r_s\, r\, a}{\Sigma^2}
    \left(2\Sigma\sin\theta\cos\theta + 2a^2\sin^3\!\theta\cos\theta\right),
\end{equation}
yang menunjukkan ketergantungan eksplisit pada sudut polar $\theta$. Komponen ini juga lenyap pada limit $a \to 0$.

Perhitungan lengkap seluruh komponen Christoffel metrik Kerr menghasilkan puluhan suku yang kompleks. Dalam praktik, perhitungan ini dapat diverifikasi dan dilakukan secara efisien menggunakan pustaka komputasi simbolik seperti \texttt{EinsteinPy} \parencite{Bapat2020EinsteinPy}.


% ============================================================
\section{Persamaan Geodesik}
\label{sec:persamaan-geodesik}

Lintasan partikel uji bermassa dalam ruang-waktu melengkung dapat diturunkan dari prinsip aksi minimal. Lagrangian untuk partikel pada metrik $g_{\mu\nu}$ diberikan oleh:
\begin{equation}
    \mathcal{L} = \frac{1}{2}\, g_{\mu\nu}\, \dot{x}^\mu\, \dot{x}^\nu,
    \label{eq:lagrangian}
\end{equation}
dengan $\dot{x}^\mu \equiv dx^\mu/d\tau$. Penerapan persamaan Euler-Lagrange,
\begin{equation}
    \frac{d}{d\tau}\frac{\partial\mathcal{L}}{\partial\dot{x}^\mu}
    - \frac{\partial\mathcal{L}}{\partial x^\mu} = 0,
\end{equation}
menghasilkan persamaan geodesik (Persamaan~\ref{eq:geodesic}). Derivasi lengkap dengan substitusi eksplisit metrik Kerr disajikan pada Bab~\ref{chap:metode}.

Untuk metrik yang memiliki simetri, terdapat konstanta gerak yang menyederhanakan persamaan geodesik. Metrik Kerr tidak bergantung pada $t$ dan $\phi$, sehingga momentum konjugat yang bersesuaian merupakan konstanta gerak:
\begin{equation}
    p_t = \frac{\partial\mathcal{L}}{\partial\dot{t}} = g_{tt}\,\dot{t} + g_{t\phi}\,\dot{\phi} \equiv \frac{E}{c},
    \label{eq:const-E}
\end{equation}
\begin{equation}
    p_\phi = \frac{\partial\mathcal{L}}{\partial\dot{\phi}} = g_{\phi t}\,\dot{t} + g_{\phi\phi}\,\dot{\phi} \equiv -L_z,
    \label{eq:const-Lz}
\end{equation}
di mana $E$ adalah energi spesifik dan $L_z$ adalah komponen momentum sudut terhadap sumbu rotasi. Selain itu, normalisasi empat-kecepatan memberikan constraint:
\begin{equation}
    g_{\mu\nu}\,\dot{x}^\mu\,\dot{x}^\nu = c^2,
    \label{eq:normalisasi}
\end{equation}
untuk partikel bermassa (\textit{timelike geodesic}) \parencite{Carroll2004SpacetimeGeometry, Schutz2009FirstCourse}.


% ============================================================
\section{Transportasi Paralel Tensor Spin}
\label{sec:transportasi-paralel}

Orientasi giroskop yang bergerak bebas dalam ruang-waktu melengkung berubah menurut hukum transportasi paralel. Vektor spin $S^\mu$ yang diangkut secara paralel sepanjang worldline $x^\mu(\tau)$ memenuhi:
\begin{equation}
    \frac{DS^\mu}{D\tau} = \frac{dS^\mu}{d\tau}
    + \Gamma^\mu_{\nu\rho}\, S^\nu\, \frac{dx^\rho}{d\tau} = 0.
    \label{eq:spin-transport}
\end{equation}
Persamaan ini menyatakan bahwa vektor spin ``konstan'' terhadap koneksi ruang-waktu, namun dapat berubah orientasinya relatif terhadap sistem koordinat luar akibat kelengkungan \parencite{Misner1973Gravitation}.

Persamaan~\eqref{eq:spin-transport} dapat ditulis ulang dalam bentuk eksplisit sebagai sistem persamaan diferensial biasa (ODE) orde satu:
\begin{equation}
    \frac{dS^\mu}{d\tau} = -\Gamma^\mu_{\nu\rho}\, S^\nu\, u^\rho,
    \label{eq:spin-ode}
\end{equation}
dengan $u^\rho = dx^\rho/d\tau$ adalah empat-kecepatan yang diperoleh dari penyelesaian geodesik. Dalam notasi matriks, sistem ini menjadi:
\begin{equation}
    \frac{d}{d\tau}
    \begin{pmatrix} S^t \\ S^r \\ S^\theta \\ S^\phi \end{pmatrix}
    = -\mathbf{A}(\tau)
    \begin{pmatrix} S^t \\ S^r \\ S^\theta \\ S^\phi \end{pmatrix},
    \label{eq:spin-matrix}
\end{equation}
di mana elemen matriks $\mathbf{A}$ adalah:
\begin{equation}
    A^\mu_{\ \nu}(\tau) = \Gamma^\mu_{\nu\rho}\, u^\rho(\tau).
    \label{eq:matrix-A}
\end{equation}

Selain itu, vektor spin harus memenuhi constraint ortogonalitas terhadap empat-kecepatan:
\begin{equation}
    g_{\mu\nu}\, S^\mu\, u^\nu = 0,
    \label{eq:spin-constraint}
\end{equation}
yang menjamin bahwa spin bersifat ``ruang murni'' (\textit{spacelike}) dalam kerangka acuan partikel. Constraint ini terpenuhi secara otomatis jika kondisi awal $S^\mu(0)$ dipilih ortogonal terhadap $u^\mu(0)$, karena transportasi paralel melestarikan produk dalam $g_{\mu\nu} S^\mu u^\nu$ sepanjang geodesik \parencite{Carroll2004SpacetimeGeometry}.


% ============================================================
\section{Efek Relativistik di Sekitar Bumi}
\label{sec:efek-relativistik}

\subsection{Efek Geodesik (Presesi de Sitter)}
\label{subsec:efek-geodesik}

Efek geodesik adalah presesi orientasi giroskop yang bergerak dalam medan gravitasi akibat kelengkungan ruang-waktu, bahkan tanpa rotasi sumber. Efek ini muncul dari transportasi paralel vektor spin sepanjang geodesik pada metrik Schwarzschild. Dalam aproksimasi medan lemah dan orbit sirkular, laju presesi geodesik diberikan oleh \parencite{Will2014Confrontation}:
\begin{equation}
    \bm{\Omega}_{\text{geodetik}} = \frac{3}{2}\,\frac{GM}{c^2\,r^3}\; \mathbf{r} \times \mathbf{v},
    \label{eq:omega-geodetic}
\end{equation}
dengan $\mathbf{r}$ adalah vektor posisi dan $\mathbf{v}$ adalah kecepatan orbit. Untuk orbit sirkular dengan $v = \sqrt{GM/r}$, magnitudo laju presesi menjadi:
\begin{equation}
    |\bm{\Omega}_{\text{geodetik}}| = \frac{3}{2}\,\frac{(GM)^{3/2}}{c^2\, r^{5/2}}.
    \label{eq:omega-geodetic-circular}
\end{equation}


\subsection{Efek Frame-Dragging (Presesi Lense-Thirring)}
\label{subsec:frame-dragging}

Efek Lense-Thirring muncul karena rotasi massa sumber yang menyeret ruang-waktu di sekitarnya. Secara matematis, efek ini berasal dari komponen campuran $g_{t\phi} \neq 0$ pada metrik Kerr yang tidak ada pada metrik Schwarzschild. Bentuk vektor umum laju presesi Lense-Thirring pada jarak $\mathbf{r}$ dari sumber bermomen sudut $\mathbf{J}$ diberikan oleh \parencite{Lense1918LenseThirring}:
\begin{equation}
    \bm{\Omega}_{\text{LT}}(\mathbf{r})
    = \frac{G}{c^2\, r^3}\left[3\hat{\mathbf{r}}\,(\hat{\mathbf{r}} \cdot \mathbf{J}) - \mathbf{J}\right],
    \label{eq:omega-lt-general}
\end{equation}
dengan $\hat{\mathbf{r}} = \mathbf{r}/r$. Untuk orbit polar (relevan dengan konfigurasi GP-B), ekspresi ini dapat disederhanakan menjadi:
\begin{equation}
    |\bm{\Omega}_{\text{LT}}| \approx \frac{2G\,|\mathbf{J}|}{c^2\, r^3}.
    \label{eq:omega-lt-approx}
\end{equation}
Derivasi lengkap ekspresi ini dari metrik Kerr melalui reduksi medan lemah disajikan pada Bab~\ref{chap:metode}.

Laju presesi total yang dialami giroskop merupakan superposisi kedua kontribusi:
\begin{equation}
    \bm{\Omega} = \bm{\Omega}_{\text{geodetik}} + \bm{\Omega}_{\text{LT}}.
    \label{eq:omega-total}
\end{equation}


% ============================================================
\section{Eksperimen Gravity Probe B}
\label{sec:gravity-probe-b}

Eksperimen \textit{Gravity Probe B} (GP-B) merupakan misi satelit yang dirancang oleh Stanford University dan NASA untuk menguji dua prediksi Relativitas Umum: efek geodesik dan efek \textit{frame-dragging} \parencite{Turneaure1986GravityProbeBApproach}. Misi ini diluncurkan pada 20 April 2004 dan mengumpulkan data selama lebih dari satu tahun.

Instrumen utama GP-B terdiri dari empat giroskop kuarsa superpresisi yang mengorbit Bumi dalam lintasan polar pada ketinggian $\sim$642~km (radius orbit $r \approx 7{,}013 \times 10^6$~m). Orientasi referensi diarahkan ke bintang panduan IM Pegasi. Parameter fisik yang relevan ditampilkan pada Tabel~\ref{tab:param-gpb}.

\begin{table}[ht]
\centering
\caption{Parameter fisik yang digunakan dalam analisis GP-B.}
\label{tab:param-gpb}
\begin{tabular}{lll}
\toprule
\textbf{Parameter} & \textbf{Simbol} & \textbf{Nilai} \\
\midrule
Massa Bumi             & $M$    & $5{,}972 \times 10^{24}$~kg \\
Momentum sudut Bumi    & $J$    & $5{,}86 \times 10^{33}$~kg\,m$^2$\,s$^{-1}$ \\
Radius orbit           & $r$    & $7{,}013 \times 10^{6}$~m \\
Radius Schwarzschild   & $r_s$  & $8{,}87 \times 10^{-3}$~m \\
Parameter rotasi       & $a$    & $3{,}98 \times 10^{-3}$~m \\
Konstanta gravitasi    & $G$    & $6{,}674 \times 10^{-11}$~m$^3$\,kg$^{-1}$\,s$^{-2}$ \\
Kecepatan cahaya       & $c$    & $2{,}998 \times 10^{8}$~m\,s$^{-1}$ \\
\bottomrule
\end{tabular}
\end{table}

Hasil akhir pengukuran GP-B yang dipublikasikan oleh \textcite{Everitt2011GravityProbeB} adalah:
\begin{equation}
    \Omega_{\text{geodetik}}^{\text{GP-B}} = 6602 \pm 18~\text{mas/yr}, \qquad
    \Omega_{\text{LT}}^{\text{GP-B}} = 37{,}2 \pm 7{,}2~\text{mas/yr},
    \label{eq:gpb-results}
\end{equation}
di mana 1~mas (milliarcsecond) $= 4{,}848 \times 10^{-9}$~rad. Kedua hasil ini konsisten dengan prediksi Relativitas Umum. Konversi satuan dari rad/s ke mas/yr menggunakan faktor:
\begin{equation}
    1~\text{rad/s} = 206\,264\,806{,}247 \times 3{,}15576 \times 10^7~\text{mas/yr}
    \approx 6{,}51 \times 10^{15}~\text{mas/yr}.
    \label{eq:konversi-masyr}
\end{equation}

Selain GP-B, efek Lense-Thirring juga dikonfirmasi oleh analisis orbit satelit LAGEOS dengan ketidakpastian sekitar 10\% \parencite{CiufoliniPavlis2004_LAGEOS}. Berbagai pendekatan alternatif, termasuk model gravitasi vektor \parencite{BeheraBarik2020} dan teori modifikasi paritas \parencite{Mao2006, Nakamura2018}, telah diajukan namun belum menunjukkan deviasi signifikan dari GR pada tingkat presisi eksperimen saat ini.


% ============================================================
\section{Pustaka Komputasi Python}
\label{sec:pustaka-python}

Penyelesaian numerik persamaan geodesik dan transportasi paralel spin memerlukan implementasi komputasi yang efisien. Penelitian ini memanfaatkan pustaka Python berikut sebagai alat bantu:
\begin{enumerate}[label=(\roman*)]
    \item \textbf{EinsteinPy} \parencite{Bapat2020EinsteinPy}: pustaka sumber terbuka untuk komputasi tensorial GR, menyediakan kelas untuk mendefinisikan metrik, menghitung simbol Christoffel, dan mengintegrasikan geodesik.
    \item \textbf{FANTASY} \parencite{Christian2020FANTASY}: integrator geodesik simplektik dengan diferensiasi otomatis, menawarkan stabilitas numerik jangka panjang.
    \item \textbf{KerrGeoPy} \parencite{Park2024KerrGeoPy}: pustaka khusus untuk geodesik pada metrik Kerr.
    \item \textbf{SciPy} \parencite{SciPy2020}: pustaka komputasi ilmiah yang menyediakan solver ODE (\texttt{solve\_ivp}) dengan berbagai metode termasuk Runge-Kutta orde 4/5 (RK45).
    \item \textbf{Matplotlib} \parencite{Matplotlib2007}: pustaka visualisasi data untuk menghasilkan grafik evolusi spin dan analisis kestabilan numerik.
\end{enumerate}

Pustaka-pustaka ini digunakan untuk memverifikasi hasil derivasi analitik dan menghasilkan solusi numerik yang disajikan pada Bab~\ref{chap:hasil}. Penjelasan detail mengenai penggunaan metode numerik disajikan pada Bab~\ref{chap:metode}.


% === BAB III: MODEL MATEMATIS DAN METODE (versi ringkas sempro) ===
% !TEX root = ../mainsempro.tex

% ==================================================
% BAB III - METODOLOGI PENELITIAN (Format Sempro UNS)
% ==================================================

\chapter{Metodologi Penelitian}
\label{chap:metode}

% ============================================================
\section{Tempat dan Waktu}
\label{sec:tempat-waktu}

Pelaksanaan penelitian ini berada di Lingkungan Program Studi Fisika, Fakultas Matematika dan Ilmu Pengetahuan Alam, Universitas Sebelas Maret Surakarta. Penelitian ini diajukan dan akan dilaksanakan pada bulan Agustus 2025 hingga bulan Desember 2025.

% ============================================================
\section{Alat dan Bahan}
\label{sec:alat-bahan}

Penelitian ini merupakan penelitian kajian analitik dan komputasi numerik yang berfokus pada derivasi matematis dan penyelesaian numerik persamaan geodesik serta transportasi paralel tensor spin pada metrik Kerr. Oleh karena itu, alat dan bahan yang digunakan dalam penelitian ini tidak berupa peralatan eksperimen laboratorium, melainkan berupa perangkat komputasi serta komponen teoretis yang mendukung proses perumusan model, penyelesaian persamaan, dan analisis hasil. Uraian mengenai alat dan bahan yang digunakan dalam penelitian ini disajikan secara sistematis pada subbab berikut:

\subsection{Alat}
\label{subsec:alat}

Peralatan utama yang digunakan dalam penelitian ini adalah perangkat komputer berupa laptop yang berfungsi sebagai sarana komputasi dan pengolahan data. Laptop yang digunakan memiliki spesifikasi prosesor Intel® Core™ i5 (quad-core) dengan kapasitas memori sebesar 16 GB serta penyimpanan 512 GB SSD.

Perangkat lunak komputasi yang digunakan dalam penelitian ini adalah Python 3.12, yang dimanfaatkan untuk melakukan perhitungan numerik, pemrosesan data, serta visualisasi hasil perhitungan. Pustaka Python utama yang digunakan meliputi EinsteinPy untuk komputasi tensorial Relativitas Umum, SciPy untuk penyelesaian sistem persamaan diferensial biasa menggunakan metode Runge-Kutta, serta Matplotlib untuk visualisasi data dan hasil simulasi. Selain itu, pustaka KerrGeoPy dan FANTASY digunakan sebagai alternatif solver geodesik pada metrik Kerr.

\subsection{Bahan}
\label{subsec:bahan}

Bahan kajian utama dalam penelitian ini adalah model ruang-waktu Kerr yang mendeskripsikan geometri ruang-waktu di sekitar massa berotasi. Geometri ruang-waktu dinyatakan melalui metrik statis dan aksial-simetris dalam koordinat Boyer-Lindquist yang melibatkan parameter massa $M$ dan parameter rotasi $a = J/(Mc)$. Pemilihan parameter orbit disesuaikan dengan konfigurasi eksperimen \textit{Gravity Probe B}, yaitu orbit polar sirkular pada ketinggian $\sim$642 km di atas permukaan Bumi.

Model fisis yang digunakan dalam penelitian ini adalah dinamika giroskop yang dideskripsikan oleh vektor spin $S^\mu$ yang mengalami transportasi paralel sepanjang lintasan geodesik. Persamaan geodesik diturunkan dari Lagrangian metrik Kerr melalui metode Euler-Lagrange, sedangkan persamaan transportasi paralel diformulasikan dalam bentuk sistem matriks $4 \times 4$.

Selain model geometri dan persamaan gerak, bahan penelitian ini juga mencakup formulasi laju presesi dalam aproksimasi medan lemah, yaitu presesi geodesik (de Sitter) dan presesi Lense-Thirring (\textit{frame-dragging}), sebagai kuantitas fisis yang dianalisis. Seluruh model dan formulasi yang digunakan dalam penelitian ini disusun berdasarkan kerangka teoritis dan referensi ilmiah yang relevan dengan kajian metrik Kerr, efek relativistik, dan penyelesaian numerik persamaan diferensial.

% ============================================================
\section{Prosedur Penelitian}
\label{sec:prosedur}

Prosedur penelitian dalam skripsi ini disusun secara sistematis untuk menggambarkan alur kerja penelitian yang mencakup tahapan analisis teoretis dan komputasi numerik. Secara umum, penelitian ini dibagi ke dalam tiga tahap utama, yaitu tahap derivasi analitik persamaan gerak, tahap implementasi dan penyelesaian numerik, serta tahap validasi hasil terhadap data eksperimen. Setiap tahap dirancang untuk menunjukkan keterkaitan yang jelas antara perumusan model fisik, penerapan metode matematis, serta perolehan hasil numerik yang dianalisis secara fisis.

Pada Tahap 1 dilakukan analisis teoretis yang mencakup penulisan eksplisit metrik Kerr, derivasi persamaan geodesik dari Lagrangian, dan formulasi persamaan transportasi paralel tensor spin dalam bentuk matriks. Tahap 2 merupakan tahapan implementasi numerik yang mencakup reduksi sistem persamaan ke orde satu dan penyelesaian menggunakan metode Runge-Kutta. Tahap 3 merupakan tahapan validasi dan analisis hasil terhadap data eksperimen \textit{Gravity Probe B}.

\subsection{Tahap 1: Derivasi Analitik Persamaan Gerak pada Metrik Kerr}
\label{subsec:tahap1}

Tahapan ini dilakukan melalui beberapa prosedur penelitian yang saling berkaitan. Penjelasan rinci mengenai setiap prosedur disajikan pada uraian berikut.

\begin{enumerate}[label=\arabic*)]
    \item Metrik Kerr dalam koordinat Boyer-Lindquist

    Menetapkan geometri Kerr sebagai latar ruang-waktu sistem fisis, dengan elemen garis yang memuat parameter massa $M$ dan parameter rotasi $a = J/(Mc)$.

    \item Lagrangian metrik Kerr

    Menyusun Lagrangian partikel bermassa $\mathcal{L} = \tfrac{1}{2}\,g_{\mu\nu}\,\dot{x}^\mu\,\dot{x}^\nu$ secara eksplisit dari komponen-komponen metrik Kerr.

    \item Konstanta gerak

    Mengidentifikasi koordinat siklik ($t$ dan $\phi$) dan menurunkan dua konstanta gerak: energi spesifik $E$ dan momentum sudut $L_z$.

    \item Persamaan geodesik

    Menerapkan persamaan Euler-Lagrange untuk koordinat $r$ dan $\theta$ guna memperoleh persamaan gerak orde dua.

    \item Transportasi paralel tensor spin

    Memformulasikan persamaan transportasi paralel $dS^\mu/d\tau = -\Gamma^\mu_{\nu\rho}\, S^\nu\, u^\rho$ dalam bentuk sistem matriks $4 \times 4$.

    \item Reduksi medan lemah

    Melakukan ekspansi metrik Kerr pada limit $r_s/r \ll 1$ dan $a/r \ll 1$ untuk memperoleh metrik terlinearisasi dan ekspresi analitik laju presesi.

    \item Ekspresi laju presesi

    Memperoleh ekspresi presesi geodesik $\bm{\Omega}_{\text{geodetik}} = \frac{3}{2}\frac{GM}{c^2 r^3}\,\mathbf{r}\times\mathbf{v}$ dan presesi Lense-Thirring $\bm{\Omega}_{\text{LT}} = \frac{G}{c^2 r^3}[3\hat{\mathbf{r}}(\hat{\mathbf{r}}\cdot\mathbf{J}) - \mathbf{J}]$.

    \item Analisis

    Memverifikasi konsistensi dimensi dan limit khusus ($a \to 0$ mereduksi ke Schwarzschild).
\end{enumerate}

\subsection{Tahap 2: Implementasi dan Penyelesaian Numerik}
\label{subsec:tahap2}

Tahapan ini dilakukan melalui beberapa prosedur penelitian yang saling berkaitan. Penjelasan rinci mengenai setiap prosedur disajikan pada uraian berikut.

\begin{enumerate}[label=\arabic*)]
    \item Reduksi ke sistem ODE orde satu

    Mendefinisikan variabel bantu $u^\mu \equiv \dot{x}^\mu$ untuk mengubah persamaan geodesik orde dua menjadi sistem orde satu, menghasilkan total 12 ODE (4 posisi + 4 kecepatan + 4 spin).

    \item Penetapan kondisi awal

    Menetapkan kondisi awal orbit sirkular polar pada radius $r_0 = R_\oplus + 642$ km dengan kecepatan Keplerian $v_0 = \sqrt{GM/r_0}$, serta vektor spin awal yang ortogonal terhadap empat-kecepatan.

    \item Perhitungan simbol Christoffel

    Menghitung seluruh komponen simbol Christoffel $\Gamma^\mu_{\nu\rho}$ dari metrik Kerr pada setiap langkah integrasi, menggunakan pustaka EinsteinPy.

    \item Integrasi numerik

    Menyelesaikan sistem 12 ODE secara simultan menggunakan metode Runge-Kutta orde empat (RK4) atau metode adaptif RK45 melalui \texttt{scipy.integrate.solve\_ivp}.

    \item Monitoring kestabilan

    Memantau tiga besaran konservatif pada setiap langkah: normalisasi empat-kecepatan ($g_{\mu\nu}u^\mu u^\nu = -c^2$), ortogonalitas spin-kecepatan ($g_{\mu\nu}S^\mu u^\nu = 0$), dan norma spin ($g_{\mu\nu}S^\mu S^\nu = \text{konstan}$).

    \item Penyimpanan data

    Menyimpan hasil integrasi dalam format numerik (\texttt{.csv}) yang memuat waktu proper $\tau$, posisi $x^\mu$, kecepatan $u^\mu$, dan komponen spin $S^\mu$.
\end{enumerate}

\subsection{Tahap 3: Validasi dan Analisis Hasil}
\label{subsec:tahap3}

Tahapan ini dilakukan melalui beberapa prosedur penelitian yang saling berkaitan. Penjelasan rinci mengenai setiap prosedur disajikan pada uraian berikut.

\begin{enumerate}[label=\arabic*)]
    \item Ekstraksi laju presesi

    Mengekstraksi sudut presesi total dari evolusi komponen spin $S^\mu(\tau)$ dengan membandingkan orientasi akhir giroskop terhadap orientasi awalnya.

    \item Pemisahan efek presesi

    Menjalankan simulasi pertama dengan metrik Schwarzschild ($a = 0$) untuk memperoleh presesi geodesik, dan simulasi kedua dengan metrik Kerr ($a \neq 0$) untuk memperoleh presesi total. Efek \textit{frame-dragging} dihitung dari selisih keduanya.

    \item Konversi satuan

    Mengkonversi laju presesi dari rad/s ke milliarcsecond per year (mas/yr) menggunakan faktor konversi $1~\text{rad/s} \approx 6{,}51 \times 10^{15}~\text{mas/yr}$.

    \item Perbandingan dengan data GP-B

    Membandingkan hasil simulasi dengan nilai eksperimen: $\Omega_{\text{geodetik}} = 6602 \pm 18$ mas/yr dan $\Omega_{\text{LT}} = 37{,}2 \pm 7{,}2$ mas/yr \parencite{Everitt2011GravityProbeB}.

    \item Perhitungan deviasi

    Menghitung tingkat kesalahan relatif $\delta = |\Omega_{\text{simulasi}} - \Omega_{\text{teori}}|/\Omega_{\text{teori}} \times 100\%$.

    \item Analisis sensitivitas

    Mengkaji pengaruh variasi parameter rotasi $a$, langkah waktu $\Delta\tau$, dan radius orbit $r$ terhadap hasil simulasi.

    \item Analisis limit khusus

    Memverifikasi bahwa pada limit $a \to 0$, hasil simulasi mereduksi ke prediksi metrik Schwarzschild, sebagai langkah validasi konsistensi model.
\end{enumerate}

% ============================================================
\section{Teknik Analisis Data}
\label{sec:teknik-analisis}

Secara umum, teknik analisis data dalam penelitian ini dilakukan secara deskriptif dan analitis untuk menginterpretasikan hasil perhitungan laju presesi geodesik dan \textit{frame-dragging} dari dinamika giroskop dalam latar ruang-waktu Kerr. Hasil analisis disajikan dalam bentuk nilai numerik, tabel, dan grafik, yang kemudian diinterpretasikan untuk menjelaskan pengaruh parameter geometri dan rotasi massa terhadap presesi giroskop.

Analisis presesi giroskop dilakukan berdasarkan evolusi temporal vektor spin $S^\mu(\tau)$ yang diperoleh dari penyelesaian numerik persamaan transportasi paralel. Laju presesi dihitung dari perubahan orientasi spin per satuan waktu dan dikonversi ke satuan mas/yr. Pemisahan kontribusi geodesik dan Lense-Thirring dilakukan melalui perbandingan hasil simulasi pada metrik Schwarzschild ($a = 0$) dan metrik Kerr ($a \neq 0$). Pengaruh variasi parameter orbit dan parameter rotasi Bumi terhadap karakteristik presesi dianalisis secara kuantitatif melalui tren numerik yang diperoleh.

Selanjutnya, analisis kestabilan numerik dilakukan dengan memantau tiga besaran konservatif sepanjang integrasi: normalisasi empat-kecepatan, ortogonalitas spin terhadap kecepatan, dan norma spin. Deviasi besaran-besaran tersebut dari nilai teoritisnya digunakan sebagai indikator akurasi dan kestabilan metode integrasi yang diterapkan. Analisis konvergensi juga dilakukan dengan memvariasikan langkah waktu $\Delta\tau$ untuk memastikan bahwa hasil simulasi tidak sensitif terhadap pilihan ukuran langkah.

Sebagai langkah validasi, dilakukan perbandingan langsung antara hasil simulasi numerik dengan prediksi analitik medan lemah serta data eksperimen \textit{Gravity Probe B}. Analisis limit khusus, seperti limit tanpa rotasi ($a \to 0$), juga dilakukan dengan membandingkan hasil yang diperoleh terhadap prediksi metrik Schwarzschild. Analisis ini bertujuan untuk memastikan konsistensi dan keabsahan fisis dari hasil penelitian.


\printbibliography

\end{document}
